\section{Results}
\label{sec:results}

We evaluate MAYA across three benchmarks (ImageNet-R, TinyImageNet, ObjectNet) and three frozen encoders (ResNet50, DINOv3, SigLIP2). We organise the results around three questions. First, does explicit geometric correction close the gap between a centroid classifier and gradient-trained replay (Section~\ref{sec:mechanics})? Second, given that MAYA stores sample-derived vectors, is that memory well spent relative to storing the vectors directly (Section~\ref{sec:matched})? Third, how does MAYA compare to analytic methods that store only sufficient statistics (Section~\ref{sec:sota})? We close by stating explicitly where MAYA is the wrong choice (Section~\ref{sec:whennot}).

All results in this section use a single seed (42) and are therefore preliminary with respect to differences below roughly two AIA points; see Section~\ref{sec:limitations}.

\subsection{Mechanics of Geometric Correction}
\label{sec:mechanics}

Table~\ref{tab:mechanics} compares MAYA against the standard replay paradigms.

% Generated by tools/make_paper_tables.py -- do not edit by hand.
\begin{table}[t]
\centering
\caption{Mechanics of geometric correction (AIA). Best result per row in bold.}
\label{tab:mechanics}
\begin{tabular}{ll|cccc}
\toprule
\textbf{Dataset} & \textbf{Backbone} & \textbf{NCM} & \textbf{ER+MLP} & \textbf{Nodes Only} & \textbf{MAYA Full} \\
\midrule
\multirow{3}{*}{ImageNet-R} & ResNet50 & 40.7\% & 46.8\% & 45.8\% & \textbf{47.6\%} \\
 & DINOv3 & 85.7\% & \textbf{94.5\%} & 93.5\% & 94.2\% \\
 & SigLIP2 & 95.2\% & \textbf{96.1\%} & 92.6\% & 95.6\% \\
\midrule
\multirow{3}{*}{TinyImageNet} & ResNet50 & 63.3\% & \textbf{71.1\%} & 64.5\% & 69.0\% \\
 & DINOv3 & 89.5\% & \textbf{94.5\%} & 90.2\% & 93.6\% \\
 & SigLIP2 & 86.1\% & \textbf{90.2\%} & 82.7\% & 89.0\% \\
\midrule
\multirow{3}{*}{ObjectNet} & ResNet50 & 19.4\% & \textbf{23.4\%} & 18.2\% & 20.0\% \\
 & DINOv3 & 47.9\% & 73.0\% & 71.1\% & \textbf{76.0\%} \\
 & SigLIP2 & 76.6\% & 80.2\% & 76.6\% & \textbf{80.5\%} \\
\bottomrule
\end{tabular}
\end{table}


\textbf{Correcting geometric bias.} The limitations of Nearest Class Mean are clearest where the frozen representation is poorly matched to the class stream. On ObjectNet with DINOv3, NCM reaches only $47.9\%$; MAYA reaches $76.0\%$. On this configuration MAYA also exceeds the gradient-trained ER+MLP head ($73.0\%$) while using less memory, without any gradient descent.

\textbf{Relative to replay.} The broader pattern is more mixed and we state it plainly: MAYA exceeds ER+MLP in three of nine configurations (ObjectNet/DINOv3, ObjectNet/SigLIP2, ImageNet-R/ResNet50) and trails it in the remaining six, by between $0.3$ and $3.4$ points. What is consistent is the memory: MAYA uses less than ER+MLP in eight of nine configurations, and on TinyImageNet it uses $38$--$44\%$ as much. The claim we make is therefore about the trade-off, not about dominance --- MAYA reaches comparable accuracy to a gradient-trained replay head at a substantially smaller and explicitly bounded memory budget.

\subsection{Is the Memory Well Spent? A Matched-Budget Comparison}
\label{sec:matched}

MAYA's episodic memory stores vectors derived from individual samples, which places it in the same family as replay methods rather than with the statistics-only analytic methods. The relevant question for such a method is not whether it uses memory, but whether it uses memory better than storing the samples directly. Table~\ref{tab:matched} answers this by giving a raw-feature exemplar buffer the same budget MAYA consumes.

% Generated by tools/make_paper_tables.py -- do not edit by hand.
\begin{table}[t]
\centering
\caption{Matched-memory comparison against raw stored feature vectors. AIA with memory footprint in MB in parentheses. At an equivalent budget, structured episodic nodes outperform an unstructured feature buffer in every configuration.}
\label{tab:matched}
\begin{tabular}{ll|ccc}
\toprule
\textbf{Dataset} & \textbf{Backbone} & \textbf{Raw vectors} & \textbf{MAYA} & \textbf{$\Delta$} \\
\midrule
\multirow{3}{*}{ImageNet-R} & ResNet50 & 37.6\% \tiny{(138)} & \textbf{47.6\%} \tiny{(160)} & $+10.0$ \\
 & DINOv3 & 69.8\% \tiny{(307)} & \textbf{94.2\%} \tiny{(383)} & $+24.4$ \\
 & SigLIP2 & 78.9\% \tiny{(100)} & \textbf{95.6\%} \tiny{(114)} & $+16.7$ \\
\midrule
\multirow{3}{*}{TinyImageNet} & ResNet50 & 56.7\% \tiny{(245)} & \textbf{69.0\%} \tiny{(246)} & $+12.3$ \\
 & DINOv3 & 88.3\% \tiny{(553)} & \textbf{93.6\%} \tiny{(555)} & $+5.3$ \\
 & SigLIP2 & 83.2\% \tiny{(178)} & \textbf{89.0\%} \tiny{(179)} & $+5.8$ \\
\midrule
\multirow{3}{*}{ObjectNet} & ResNet50 & 12.2\% \tiny{(229)} & \textbf{20.0\%} \tiny{(243)} & $+7.8$ \\
 & DINOv3 & 28.8\% \tiny{(508)} & \textbf{76.0\%} \tiny{(550)} & $+47.2$ \\
 & SigLIP2 & 66.0\% \tiny{(168)} & \textbf{80.5\%} \tiny{(177)} & $+14.5$ \\
\bottomrule
\end{tabular}
\end{table}


At a matched budget, structured episodic nodes outperform an unstructured feature buffer in \emph{all nine} configurations, by a mean of $16.0$ points and by as much as $47.2$ on ObjectNet with DINOv3. This is the clearest evidence we have that the local memory is doing structural work rather than merely caching data: the same bytes, allocated to nodes maintained under the global alignment, are worth substantially more than the same bytes allocated to stored feature vectors.

We note one caveat. The raw-vector baseline selects exemplars without a selection strategy; a herding-based selector would be a stronger comparison, and we regard the magnitude of this gap as an upper bound on the advantage.

\subsection{Comparison with Analytic Methods}
\label{sec:sota}

Table~\ref{tab:sota} compares MAYA with frozen-backbone analytic CIL methods that store only sufficient statistics.

% Generated by tools/make_paper_tables.py -- do not edit by hand.
\begin{table}[t]
\centering
\caption{Comparison with analytic continual-learning methods. AIA with memory footprint in MB in parentheses. Best accuracy per row in bold.}
\label{tab:sota}
\begin{tabular}{ll|cccc}
\toprule
\textbf{Dataset} & \textbf{Backbone} & \textbf{MAYA} & \textbf{Deep SLDA} & \textbf{Exact ACL} & \textbf{RanPAC ($M{=}10$k)} \\
\midrule
\multirow{3}{*}{ImageNet-R} & ResNet50 & 47.6\% \tiny{(160)} & 51.9\% \tiny{(18)} & 52.8\% \tiny{(18)} & \textbf{55.6\%} \tiny{(467)} \\
 & DINOv3 & 94.2\% \tiny{(383)} & 94.0\% \tiny{(67)} & 95.7\% \tiny{(67)} & \textbf{95.9\%} \tiny{(545)} \\
 & SigLIP2 & 95.6\% \tiny{(114)} & \textbf{96.6\%} \tiny{(10)} & 95.2\% \tiny{(10)} & 96.6\% \tiny{(448)} \\
\midrule
\multirow{3}{*}{TinyImageNet} & ResNet50 & 69.0\% \tiny{(246)} & 69.6\% \tiny{(18)} & 70.8\% \tiny{(18)} & \textbf{73.2\%} \tiny{(467)} \\
 & DINOv3 & 93.6\% \tiny{(555)} & 92.4\% \tiny{(67)} & 94.9\% \tiny{(67)} & \textbf{94.9\%} \tiny{(545)} \\
 & SigLIP2 & 89.0\% \tiny{(179)} & 88.2\% \tiny{(10)} & 89.0\% \tiny{(10)} & \textbf{90.6\%} \tiny{(448)} \\
\midrule
\multirow{3}{*}{ObjectNet} & ResNet50 & 20.0\% \tiny{(243)} & \textbf{24.7\%} \tiny{(18)} & 23.6\% \tiny{(18)} & 24.3\% \tiny{(472)} \\
 & DINOv3 & \textbf{76.0\%} \tiny{(550)} & 73.7\% \tiny{(69)} & 75.7\% \tiny{(69)} & 75.4\% \tiny{(550)} \\
 & SigLIP2 & 80.5\% \tiny{(177)} & 81.5\% \tiny{(11)} & 78.7\% \tiny{(11)} & \textbf{82.0\%} \tiny{(452)} \\
\bottomrule
\end{tabular}
\\[2pt]
\footnotesize Single seed (42). RanPAC's footprint is dominated by its $O(M^2)$ precision matrix; SLDA and ACL by their $O(d^2)$ covariance.
\end{table}


MAYA achieves the best accuracy of the four methods on ObjectNet with DINOv3 ($76.0\%$ against $75.7\%$ for Exact ACL), the configuration with the most out-of-distribution class geometry, and it matches Exact ACL on TinyImageNet/SigLIP2. In the remaining seven configurations an analytic baseline is more accurate. We do not claim state-of-the-art accuracy.

The memory comparison is unfavourable and we state it directly. SLDA and ACL maintain an $O(d^2)$ covariance --- $10$--$69$MB across our backbones --- while MAYA's footprint ranges from $114$ to $555$MB, since it retains that same order of statistics \emph{plus} the node bank. Where MAYA is competitive on accuracy, it is not competitive on accuracy per megabyte.

What MAYA provides that these methods structurally cannot is an explicit, addressable memory. SLDA and ACL accumulate their evidence into a single monolithic matrix from which no individual concept can be located or removed; MAYA's nodes are discrete and per-class, which permits inspection and targeted deletion. They also give up the privacy property that motivates analytic CL: retaining sample-derived vectors is a real cost, and we count it as one.

\subsection{When MAYA Is Not the Right Choice}
\label{sec:whennot}

Because accuracy-per-byte is the axis on which MAYA is weakest, we characterise it explicitly rather than leave it to be inferred. Figure~\ref{fig:pareto_audit} audits, for every configuration, whether any baseline achieves higher accuracy at or below MAYA's own memory budget.

\begin{figure}[htbp]
    \centering
    \includegraphics[width=\textwidth]{figures/pareto_audit.pdf}
    \caption{\textbf{Accuracy-memory position.} \textbf{(a)} ObjectNet/DINOv3. Sweeping $K$ traces MAYA along a curve --- its footprint is a tunable knob, which the analytic baselines do not have --- but Exact ACL sits above that curve at a fraction of the memory. \textbf{(b)} Across all nine configurations, a cheaper baseline is more accurate in eight.}
    \label{fig:pareto_audit}
\end{figure}

MAYA sits on the accuracy-memory Pareto frontier in one of nine configurations. In the other eight, a statistics-only method using between $5\times$ and $25\times$ less memory is equally or more accurate. Panel (a) shows why sweeping $K$ does not fix this: MAYA's cheapest operating point still costs more than ACL's fixed footprint while scoring far lower.

The practical reading is that MAYA is not the right choice when accuracy per megabyte is the binding constraint, and it is a reasonable choice when the deployment needs a bounded, inspectable, editable memory --- or when the frozen representation is badly matched to the class stream, which is where its accuracy advantage concentrates.
